% Generated draft: requires mathematical and source review.
\documentclass[11pt]{article}
\usepackage[margin=25mm]{geometry}
\usepackage{fontspec}
\setmainfont{Latin Modern Roman}
\usepackage{amsmath,amssymb,mathtools,unicode-math}
\setmathfont{Latin Modern Math}
\usepackage{enumitem,needspace,xurl,hyperref}
\setlength{\parindent}{0pt}
\setlength{\parskip}{5pt}
\setlength{\emergencystretch}{3em}
\newcommand{\N}{\mathbb{N}}
\newcommand{\Z}{\mathbb{Z}}
\newcommand{\Q}{\mathbb{Q}}
\newcommand{\R}{\mathbb{R}}
\newcommand{\C}{\mathbb{C}}
\newcommand{\F}{\mathbb{F}}
\newcommand{\A}{\mathbb{A}}
\newcommand{\PP}{\mathbb{P}}
\newcommand{\E}{\mathbb{E}}
\newcommand{\eps}{\varepsilon}
\newcommand{\dd}{\,\mathrm{d}}
\begin{document}
\section*{UnsolvedMath ID: 30003423}
\textbf{Source problem code:} catalogue entry
\subsection{Definitions and mathematical statement}
Let \(k\) denote the ground field. The supplied record does not specify whether \(k\) is algebraically closed or impose a characteristic hypothesis; these remain unresolved hypotheses of the intended formulation. Let \(A\) be a tame finite-dimensional \(k\)-algebra, where tameness is understood in the representation-theoretic sense of classification of indecomposable modules by finitely many one-parameter families in each dimension, apart from finitely many exceptions. The field assumptions for this usage require confirmation.

Write \(\operatorname{rep}(A,d)\) for the parameter space of representations of \(A\) with dimension datum \(d\), and \(G(d)\) for its change-of-basis group. The record does not specify whether \(d\) denotes total dimension or a dimension vector. In the intended geometric invariant theory construction, let \(\theta\) be an integral stability weight with \(\theta(d)=0\). Use the convention that a representation \(M\) is \(\theta\)-semistable if \(\theta(\dim M)=0\) and \(\theta(\dim N)\leq 0\) for every subrepresentation \(N\subseteq M\); it is \(\theta\)-stable if the inequality is strict for every nonzero proper subrepresentation. This fixes a sign convention, which the record does not specify.

For a closed \(G(d)\)-invariant subvariety \(C\subseteq\operatorname{rep}(A,d)\), write \(\mathcal M(C)_{\theta}^{ss}\) for the intended geometric invariant theory moduli space of its \(\theta\)-semistable representations. Its geometric points represent S-equivalence classes, where S-equivalence means equality of the isomorphism classes of the associated polystable representations. For a variety \(X\) and positive integer \(m\), write \(S^m(X)=X^m/\mathfrak S_m\) for its symmetric power.

\textbf{Decomposition question.} Let \(A\) be tame, and let \(C\subseteq\operatorname{rep}(A,d)\) be closed, irreducible, and invariant under \(G(d)\). Suppose that \(C\) has the \(\theta\)-stable decomposition
\[
C=m_1C_1\dotplus\cdots\dotplus m_rC_r,
\qquad r\geq 1,\quad m_i\in\Z_{>0},
\]
in the sense used by the source: the stable factors of a general semistable representation are grouped into the families \(C_i\) with multiplicities \(m_i\). The notation refers to stable factors and is not an assertion that every representation in \(C\) is a direct sum. Suppose that no \(C_i\) is the closure of a single change-of-basis orbit. The record presents the direct-sum construction as a natural finite birational surjective morphism
\[
\Psi:\prod_{i=1}^r S^{m_i}\!\left(\mathcal M(C_i)_{\theta}^{ss}\right)
\longrightarrow \mathcal M(C)_{\theta}^{ss}.
\]
Is \(\Psi\) an isomorphism for every such choice of data? Its stated existence and properties are retained as supplied premises; any additional hypotheses needed by the underlying decomposition theorem require source review.

\textbf{Smoothness question.} Are moduli spaces of semistable representations of tame finite-dimensional algebras always smooth? Here smoothness is intended over the ground field \(k\). The record does not determine whether this quantifies over quotients of entire representation spaces, irreducible components, or all closed irreducible invariant subvarieties. Consequently, the precise scope of this second universal question remains unresolved.

The supplied dated assessment reports an isomorphism result under normality and positive results in tame and special-biserial settings, while describing the general questions as open. It does not specify the exact normality hypothesis or the scope of those positive results; these statements are retained as curation claims rather than independently established conclusions.
\subsection{Short English statement}
For a tame finite-dimensional algebra, the first question asks whether the supplied finite birational map from symmetric powers of stable-factor moduli spaces to the original moduli space is always an isomorphism. The second asks whether semistable representation moduli spaces of tame algebras are always smooth. Ground-field hypotheses, the exact scope of the smoothness question, and the decomposition theorem's applicability require review.
\subsection{Sources}
\begin{itemize}
\item Local registry: \texttt{data/unsolvedmath/problems.json}, supplied record with permanent UnsolvedMath ID \(30003423\), titled \emph{Smooth Moduli Spaces of Tame Algebras}; exact source code \texttt{catalogue entry}. No individual proposer is supplied. The proposed year is recorded as 2017.
\item Supplied report citation: \emph{Representation Theory of Quivers and Finite Dimensional Algebras (2018)}, Oberwolfach Reports, pp.~591--681, DOI \texttt{10.4171/owr/2017/12}. \url{https://doi.org/10.4171/owr/2017/12}. The citation's 2018 label and its 2017 DOI/proposed-year information are retained without reconciliation.
\item Supplied literature source: \emph{Decomposing moduli of representations}. \url{https://arxiv.org/abs/1705.10255}.
\item Supplied literature source: \emph{Invariant theory for tame tilted algebras}. \url{https://arxiv.org/abs/1109.2915}.
\item Supplied literature source: \emph{OWR Representation Theory report}. \url{https://ems.press/content/serial-article-files/46676}.
\item Literature assessment dated exactly \texttt{2026-08-21}: the supplied curation labels the problem partially solved, reports an isomorphism result under normality and positive results for tame and special-biserial cases, and describes the all-tame and universal smoothness questions as open. This is supplied dated curation, not an independently verified literature assessment.
\item The registry labels the statement \texttt{corrected\_verified} and reports that it was checked against the official Oberwolfach report. That verification is a supplied provenance claim; no source was independently consulted for this draft.
\end{itemize}
\end{document}
